\chapter*{Abstract}
\addcontentsline{toc}{chapter}{Abstract}

This thesis investigates whether retrieval-augmented generation and parameter-efficient fine-tuning improve structured dashboard-design recommendations produced by instruction-tuned language models. A provenance-aware dataset was constructed from source-grounded analytical records and constrained, explicitly marked LLM-generated enrichment. Four model profiles were evaluated under prompt-only, retrieval-augmented, QLoRA-fine-tuned, and combined conditions. The executed matrix contains 32 complete runs on 274 held-out items: three seeds for Qwen3-1.7B and OLMo-2-1B-Instruct, and one seed for Qwen3-8B and Qwen3.8-27B.

The evaluation separates JSON recovery, runtime-contract validity, completeness, reference agreement at three scoring layers, per-class performance, robustness, grounding, and latency. This decomposition is necessary because end-to-end chart accuracy confounds the chart decision with output formatting and truncation. At the text level, fine-tuning improves agreement with the frozen chart reference in all eight model-by-seed blocks by 0.36--10.58 percentage points, although output validity falls for Qwen3-8B and OLMo under the executed 512-token generation budget. Retrieval is slightly beneficial for Qwen3-1.7B, improves recoverability but not chart choice for OLMo, and significantly lowers reference agreement for the 8B and 27B Qwen3 models. A traceable conflict between the knowledge base and the frozen reference explains most pie-to-donut substitutions at both larger scales. Adding retrieval after fine-tuning lowers reference agreement in eleven of twelve blocks.

The automatic results are conditional on an internally constructed reference, asymmetric seed coverage, heterogeneous code states, a small three-document knowledge base, and text-only recommendations. A completed exploratory human evaluation of the four Qwen3.8-27B conditions contributed 81 rating pages from ten complete sessions across eight briefs. On a five-point rubric, retrieval-only achieved the highest equal-brief means for chart appropriateness (4.48), rationale quality (4.52), and overall usefulness (4.52); prompt-only was marginally higher for layout and interaction. Fine-tuned conditions did not improve the human-rated outcomes in this sample. Because recruitment was by convenience, coverage was two to four pages per output, and ordinal inter-rater agreement was low for five of six dimensions, these observations are descriptive rather than population-level or causal conclusions. The thesis contributes a reproducible, artifact-anchored evaluation of structured recommendation generation and demonstrates why format conformance, semantic agreement, and human judgement must be reported as separate evidence layers.

\textbf{Keywords:} large language models; dashboard recommendation; QLoRA; retrieval-augmented generation; structured generation; provenance; evaluation
